\section{Introduction}
An agent that reads an order through a web API can present an access token to the
resource server. An agent that calls a business instead presents sound to an IVR
or a person. The latter interface exposes information and accepts requests before
the receiving organization necessarily knows which software key is calling,
which customer it represents, or what that customer authorized. A plausible
voice and a plausible request do not answer those questions.

This problem is visible in deployed contact centers. In an October 9, 2026 field
report, Pindrop describes detecting Meta Muse calls by matching two synthetic
voices, Brett and Hailey. It reports eightfold weekly growth between the first
and fourth observed weeks, in a September 9--October 6 customer sample, and notes
that different customers classify such calls differently. The report explicitly
distinguishes recognizing an agent from establishing its customer's identity and
permission~\cite{pindrop2026}. This is vendor-reported observational evidence of
an emerging operational problem, not a representative population measurement or
an independent validation of detection accuracy. We use it to motivate the
authorization question, not to estimate fraud prevalence or justify a claim
about all personal agents.

The broader issue concerns where authority enters a multimodal system. A model
may receive text, audio, images, and sensor observations, while holding access to
tools with real effects. Research on language-conditioned robotics illustrates
how linguistic and visual inputs can influence physical actions~\cite{saycan,rt2}.
Earlier microphone injection attacks show why perceptual plausibility does not
constitute authorization~\cite{dolphin,lightcommands}; recent auditory prompt
injection work extends this concern to audio-language models~\cite{audiohijack}.
A transport securing one service connection does not automatically protect a
separate microphone, rendered page, or telephone leg.

We investigate a concrete design principle: \emph{when an interaction is mediated
by a modality, the receiving executor should obtain explicit authority bound to
the endpoint participating through that modality}. In voice, this can mean
carrying a cryptographic control exchange through call audio. This principle
does not require encrypting every sound, proving who physically produced it,
or replacing an available authenticated API. It requires preventing an
unauthenticated observation from silently becoming sufficient authority for an
account action.

\sys{} instantiates this principle for participating voice agents and business
gateways. It transmits protected structured requests through an audio modem,
authenticates endpoint certificates with TLS 1.3, and separately checks signed
customer permissions. The language model may propose an operation but does not
issue its own permissions or receive private keys and grant credentials. A
business handler runs only after a deterministic gateway checks the request.
Ordinary speech remains outside this protected control exchange.

Two transport profiles expose the central networking tradeoff. The original
profile carries an OpenSSL TLS byte stream over ordered audio frames with
acknowledgments and retries. TLS therefore does \emph{not} require the physical
audio path itself to be lossless. The new profile carries QUIC packets through
unreliable audio frames, using QUIC's own TLS handshake and stream recovery.
Removing the lower reliable layer avoids composing two recovery mechanisms,
but does not establish that QUIC is faster or more robust on every audio path.
Both profiles depend on successful modem decoding.

The prototype also bridges an HTTPS personal-agent session to a short-lived,
request-specific voice grant and supports fraud assessment before execution.
An authenticated caller may still be malicious; detection and permission checks
serve different purposes and can be used together.

This draft makes three contributions. First, it specifies the voice-to-executor
trust boundary and distinguishes endpoint identity, delegated permission, and
untrusted perceptual content. Second, it presents a working design that reuses
standard cryptographic transports over audio and composes them with external
authorization. Third, it provides reproducible local feasibility evidence and
a falsifiable evaluation plan. The evidence is deliberately narrow: separate
processes on one machine, controlled digital audio, simulated business effects,
and synthetic fault injection. We do not claim a new cryptographic primitive,
the first secure voice modem, a solution to general prompt injection, or a
completed deployment study.
